\chapter{Literature Review}
\label{ch:literature_review}

This chapter reviews the literature relevant to the research domain. Section~\ref{sec:lit_hallucination} introduces hallucination in RAG. Section~\ref{sec:lit_blackbox} reviews black-box and reference-based detection. Section~\ref{sec:lit_probing} reviews internal-state probing. Section~\ref{sec:lit_ffn} discusses FFN layers as factual memory. Section~\ref{sec:lit_closedloop} reviews closed-loop and self-corrective generation, and Section~\ref{sec:lit_crossarch} discusses cross-architecture generalization.

%% -------------------------------------------------
\section{Hallucination in Retrieval-Augmented Generation}
\label{sec:lit_hallucination}

Hallucination refers to the generation of fluent but factually unsupported, unverifiable, or source-contradictory content. Prior surveys identify hallucination as a central reliability problem in large language models and commonly distinguish between \emph{intrinsic hallucination}, where the output contradicts the provided source, and \emph{extrinsic hallucination}, where the output introduces plausible but unsupported information beyond the available evidence~\cite{ji2023hallucination,huang2023survey,rawte2023survey}.

%% -------------------------------------------------
\section{Black-Box, Statistical, and Reference-Based Hallucination Detection}
\label{sec:lit_blackbox}

SelfCheckGPT is a representative black-box method. It samples multiple responses from the same model and measures whether the factual content remains consistent across samples~\cite{manakul2023selfcheckgpt}. The method is attractive because it does not require access to model weights, hidden states, or external databases. However, it requires multiple generations per query and can fail when the model consistently repeats the same incorrect claim.

Reference-based methods evaluate generated claims against retrieved or provided evidence. RefChecker, for example, decomposes responses into fine-grained claim triplets and checks them against references~\cite{hu2024refchecker}. Such approaches are useful for post-hoc verification, but they usually require an additional checking pipeline and do not directly use the generating model's internal computation as the intervention signal.

Overall, black-box, statistical, and reference-based methods provide important baselines, but they are not sufficient for the closed-loop setting studied here because they either operate after generation, require repeated sampling, depend on external checkers, or do not directly control the RAG inference loop.

%% -------------------------------------------------
\section{Internal-State Probing for Hallucination Detection}
\label{sec:lit_probing}

Internal-state probing methods are motivated by the observation that transformer activations can contain factuality-relevant information not fully visible from surface text or output probabilities. Instead of judging only the final answer, these methods inspect hidden representations produced inside the model.

Azaria and Mitchell~\cite{azaria2023saplma} showed that hidden activations can predict whether statements are true or false. Their classifier, trained on internal states, achieves approximately \textbf{71--83\% accuracy} depending on the base model. This demonstrates that internal activations can contain truthfulness-related information. However, SAPLMA-style probing is primarily a detector: it produces a factuality score but does not connect that score to an inference-time RAG correction policy.

INSIDE/EigenScore extends internal-state analysis by measuring representation-level uncertainty through the eigenspectrum or covariance structure of hidden states across multiple sampled generations~\cite{chen2024inside}. This provides a useful internal-state uncertainty signal, but it still depends on multiple generations, increasing inference cost.

MIND, proposed by Su et al.~\cite{su2024mind}, is also relevant because it uses LLM internal states for real-time hallucination detection without manual annotation. This supports the broader claim that internal representations can be useful for hallucination detection during inference. However, like most detection-oriented methods, it does not implement the same multi-action RAG controller used in the present work.

Together, these methods show that hallucination detection can benefit from hidden states, attention behavior, FFN activity, uncertainty signals, and context--knowledge balance. The remaining gap is to connect such detection signals to a closed-loop inference-time controller that can change the RAG system's behavior before a risky answer reaches the user.

%% -------------------------------------------------
\section{FFN Layers as Factual Memory and Motivation for IAV Probing}
\label{sec:lit_ffn}

The motivation for probing FFN intermediate activations comes from mechanistic interpretability work on transformer feed-forward layers. Geva et al.~\cite{geva2021ffn} showed that transformer FFN layers can operate as key-value memories: keys correspond to textual patterns, while values influence the output-token distribution. They also observed that lower layers tend to capture shallower patterns, whereas upper layers capture more semantic information.

These findings motivate the Intermediate Activation Vector (IAV) used in this work. The residual-stream CEV reflects the accumulated contextual state after attention and feed-forward processing, while the IAV captures the FFN contribution before it is projected back into the residual stream. Combining CEV and IAV therefore allows the detector to examine both context-integration behavior and parametric-memory contribution.

%% -------------------------------------------------
\section{Closed-Loop and Self-Corrective Generation}
\label{sec:lit_closedloop}

Prior work on self-corrective generation usually follows two directions: model self-reflection and retrieval correction. Self-RAG trains a language model to retrieve, generate, and critique using special reflection tokens~\cite{asai2024selfrag}. This improves controllability, but it requires training the model with reflection-token supervision and is therefore different from a frozen-model activation-probing approach.

CRAG improves RAG robustness by adding a lightweight retrieval evaluator that assesses retrieved-document quality and triggers corrective retrieval actions, including web search when static retrieval is insufficient~\cite{yan2024crag}. CRAG directly addresses retrieval failure, but its intervention signal comes from a retrieval evaluator rather than from the generating LLM's own internal activations.

Abstention is also important for safe generation. Wen et al.~\cite{wen2024abstention} survey abstention in LLMs and frame refusal as a mechanism for reducing hallucination and improving safety when the model is uncertain or when a reliable answer cannot be produced. This supports the use of ABSTAIN as one branch of the proposed controller.

To the best of our knowledge, prior work has not \emph{combined hallucination detection from the generating model's own CEV/IAV activations with a multi-action inference-time RAG controller that can choose among accepting, regenerating, re-retrieving, and abstaining in a frozen-model setting.}

%% -------------------------------------------------
\section{Cross-Architecture Generalization}
\label{sec:lit_crossarch}

Many hallucination-detection methods are evaluated on one model family or a limited set of related backbones. HalluRAG~\cite{frohling2024hallurag} is relevant here because it studies closed-domain RAG hallucination using internal states and reports that MLP classifiers trained on HalluRAG can reach test accuracies up to \textbf{75\%}, while also showing limited generalizability across settings. This motivates evaluation beyond a single model or dataset.

Cross-Layer Attention Probing (CLAP) further addresses generalization by evaluating fine-grained hallucination detection across \textbf{five LLMs and three tasks}~\cite{zhao2024clap}. This supports the need to examine whether activation-based hallucination signals transfer across different model families.

The present work evaluates the same CEV/IAV probing and controller design across three decoder-only transformer backbones: Mistral-7B-Instruct-v0.2, LLaMA-3-8B-Instruct, and Qwen3-8B. These models differ in tokenizer design, training data, layer structure, and FFN intermediate dimensions. This provides a controlled test of whether the proposed activation-probing framework can be reused across related decoder-only transformer families without modifying the backbone architecture.
